\begin{abstract}
This thesis studies protein--protein binding affinity prediction from structure and asks one focused question: can short explicit-solvent molecular dynamics improve a residue-graph affinity model beyond a strong static baseline? The pipeline used in this study starts from a curated wild-type subset of PPB-Affinity, runs short molecular-dynamics simulations for each complex, constructs residue graphs with DeepRank2, adds molecular-dynamics-derived node and edge summaries, and trains either a masked graph variational autoencoder or a direct supervised graph regressor.

The final evaluation uses a matched 1437-complex subset for which both full-complex and interface-only graphs are available. The final model matrix contains 52 configurations evaluated under repeated outer resampling with 5 folds, 2 repeats, and 520 total fits. Two graph views are compared, full and interface, together with two feature modes, static ($S$) and static-plus-dynamic ($SD$).

The main results are clear. The full graph outperforms the interface graph in all 24 matched variational-model comparisons. The dynamic augmentation does not provide a robust held-out improvement over the static representation. The strongest overall model is the direct supervised baseline on the full static graph, with mean test RMSE 2.3379, Pearson correlation 0.5064, and $R^2$ 0.2146. Semi-supervised dynamic models fit the training data more strongly but show larger train--test gaps.

The conclusion is therefore narrow but firm. In this pipeline, the molecular-dynamics-derived encoding used here does not improve affinity prediction over a well-constructed static residue graph. The result identifies the main directions for further work: richer dynamic representations, longer simulations, and less restrictive interface definitions.
\end{abstract}
